\section{相关技术与理论基础}

本章梳理了支撑本课程设计的五项核心技术背景：视觉语言模型、CLIP 语义空间、EEG 视觉解码、视觉退化建模与多模态融合，以及 LoRA 参数高效微调。各小节仅聚焦与本设计密切相关的核心原理与实现方法。

\subsection{视觉语言模型概述（VLM）}

视觉语言模型（Vision-Language Model, VLM）是能够联合处理图像与文本两种模态信息，并在此基础上实现跨模态理解与生成的深度神经网络。依据输出形态的不同，现有 VLM 可归为对齐型与生成型两大类别。

对齐型 VLM 以 CLIP 为典型代表。CLIP 在大约 4 亿对图文样本上执行对比式预训练，通过 InfoNCE 损失使图像编码器（ViT 或 ResNet）与文本编码器（Transformer）学得统一的嵌入空间。训练完成后，语义内容相近的图像与文本在该空间中彼此聚拢。因此，仅靠余弦相似度比较即可完成零样本分类任务。对齐型 VLM 产出固定维度的嵌入向量，适用于检索、分类和特征提取等场景。

生成型 VLM 在对齐结构的基础上额外引入语言解码模块，代表性模型包括 BLIP-2、LLaVA 与 Qwen-VL 系列。BLIP-2 设计了 Q-Former 作为视觉与语言之间的可学习桥接：视觉特征经查询 token 压缩后送入冻结的大语言模型。LLaVA 采用了更简洁的策略——仅通过单层线性投影将视觉 token 映射至语言模型嵌入空间，再借助视觉指令微调赋予模型对话与多轮交互能力。Qwen-VL 是阿里巴巴通义千问团队开发的原生多模态系列，其视觉编码器与语言模型在预训练阶段即实现了端到端图文对齐。

两者在输出形式、训练目标和适用场景上存在本质差异。对齐型 VLM 仅有编码器，输出固定维度的向量，主攻检索、分类等无需文本生成的任务；CLIP 的零样本分类精度在 ImageNet 上已接近有监督训练的 ResNet-50。生成型 VLM 则配备自回归文本解码器，可按 token 序列逐位产生自然语言文本，适用于图像描述、视觉问答与指令跟随。BLIP-2 的核心创新在于 Q-Former 作为轻量查询桥接器，将视觉特征压缩为少量可学习 token 后再注入冻结 LLM，从而回避了重新训练大语言模型的巨额开销。LLaVA 使用单一的线性投影层替代 Q-Former，结构更为简洁。Qwen-VL 在此基础上进一步优化了视觉编码器与 LLM 的原生集成，支持任意分辨率图像输入与更长的视觉 token 序列。

\begin{table}[h]
\centering
\caption{对齐型 VLM 与生成型 VLM 对比}
\label{tab:vlm-comparison}
\begin{tabular}{lcc}
\toprule
特性 & 对齐型 VLM（如 CLIP） & 生成型 VLM（如 BLIP-2 / Qwen-VL） \\
\midrule
输出形式 & 固定维度嵌入向量 & 逐 token 文本序列 \\
代表性任务 & 零样本分类、图文检索 & 图像描述、视觉问答 \\
文本解码器 & 无 & 有（自回归 LLM） \\
参数规模 & 通常较小（约 100M） & 通常较大（约 1B--7B） \\
\bottomrule
\end{tabular}
\end{table}

本设计在语义识别环节选用对齐型 VLM——CLIP，利用其零样本分类能力在 40 类的有限空间内进行 EEG 辅助分类；在生成式描述环节选用生成型 VLM——Qwen2-VL-2B-Instruct，利用其视觉-语言对齐能力与自回归生成机制输出自由形式的图像描述。

\placeholderfigure{CLIP 与生成式视觉语言模型结构区别示意图}{clip-vlm-contrast-placeholder}{3.2cm}

\subsection{CLIP 图文语义空间}

CLIP 的训练采用对称式 InfoNCE 对比损失。设有图文对集合 $\{(x_i, t_i)\}_{i=1}^{N}$，批大小为 $N$，图像编码器输出 $f_{\mathrm{img}}(x_i)$，文本编码器输出 $f_{\mathrm{text}}(t_i)$。优化目标为最大化正确配对样本间的余弦相似度，同时最小化错误配对样本间的相似度：

\begin{equation}
  \mathcal{L} = -\frac{1}{2N}\sum_{i=1}^{N}\left[
    \log\frac{\exp(\tau^{-1}\cos(f_{\mathrm{img}}(x_i), f_{\mathrm{text}}(t_i)))}
         {\sum_{j=1}^{N}\exp(\tau^{-1}\cos(f_{\mathrm{img}}(x_i), f_{\mathrm{text}}(t_j)))} +
    \log\frac{\exp(\tau^{-1}\cos(f_{\mathrm{img}}(x_i), f_{\mathrm{text}}(t_i)))}
         {\sum_{j=1}^{N}\exp(\tau^{-1}\cos(f_{\mathrm{img}}(x_j), f_{\mathrm{text}}(t_i)))}
  \right]
\end{equation}

其中 $\tau$ 是可训练的温度参数，用于调控相似度分布的集中程度。

零样本分类是 CLIP 的一项重要能力。对于 $K$ 个候选类别，先为每一类 $c$ 构造文本提示（例如 ``a photo of a \{class\_name\}''），经由 CLIP 文本编码器得到文本原型 $p_c = f_{\mathrm{text}}(t_c)$。对任意待分类图像 $x$，其预测类别即为与图像嵌入余弦相似度最高的原型所对应的类别：

\begin{equation}
  \hat{c} = \arg\max_{c\in\{1,\dots,K\}} \frac{f_{\mathrm{img}}(x)^{\top} p_c}{\|f_{\mathrm{img}}(x)\|\,\|p_c\|}
\end{equation}

我们充分利用了这一特性：预先通过 CLIP 文本编码器计算 40 个类别的文本原型，在训练与评估过程中保持这些原型冻结不变；模型只需学习将融合后的多模态嵌入推向正确的原型方向，从而避免在有限的 EEG 样本上从头训练分类器。

\subsection{EEG 视觉语义解码}

脑电图（EEG）通过头皮表面部署的电极阵列记录大脑皮层的电生理活动。在视觉认知实验中，EEG 能够捕捉被试观看图像刺激时产生的诱发响应，包括早期的 P100/N170 视觉成分和晚期的 P300 等与语义加工相关的成分。EEG 的核心优势在于毫秒级的时间分辨率，使其特别适合追踪视觉刺激诱发的快速神经编码过程。

然而，将 EEG 用于视觉语义解码面临着若干固有困难：(1) 信噪比偏低——单次 trial 的信号经常被肌电、眼动和环境电磁噪声所干扰；(2) 个体差异较大——不同被试的电极布局和颅骨厚度各不相同，导致响应模式呈现显著变异性；(3) 单 trial 数据量有限——典型实验条件下，每个被试在每个类别上的有效 trial 数仅为数十条。这些限制共同决定了直接从小样本 EEG 学习自由文本生成存在严重的数据瓶颈。

针对上述挑战，研究人员已进行了多方面的探索。Spampinato 等人率先将深度学习方法引入 EEG 视觉内容解码，在物体分类任务上取得了具有影响力的成果。Palazzo 等人构建了 EEG-ImageNet 大规模数据集，将 ImageNet 图像与同步采集的 EEG 进行配对。近年来，将 EEG 嵌入对齐到 CLIP 语义空间以及基于 EEG 条件生成图像等工作陆续涌现。但需要指出的是，既有研究大多在视觉完整清晰的条件下进行，视觉退化场景下 EEG 的语义辅助作用尚未得到系统性的对照检验。

本设计在数据处理层面采取了以下应对策略：在 CLIP 预训练语义空间中完成 EEG 与图像的对齐，借助 CLIP 在大规模数据上学到的语义结构来降低对 EEG 样本量的依赖；同时设置 shuffled EEG（打乱脑电）和 random EEG 两种对照条件，确保观察到的性能提升可归因于真实配对的 EEG 语义信息。

\subsection{视觉退化建模}

为系统考察模型在不同图像质量水平下的表现，我们定义了六种视觉退化场景：

\begin{table}[h]
\centering
\caption{本设计采用的六种视觉退化条件}
\label{tab:degradations}
\begin{tabular}{llc}
\toprule
退化类型 & 实现方式 & 模拟场景 \\
\midrule
clean & 原始图像，不施加任何退化 & 理想图像条件 \\
lowres16 & 缩放到 $16\times16$ 再放大回 $224\times224$ & 极低分辨率 \\
strong\_blur & 高斯模糊 $\sigma=5$ & 运动模糊或失焦 \\
strong\_noise & 高斯噪声标准差 0.15 & 传感器噪声 \\
occlusion50 & 随机遮盖 50\% 区域（Cutout） & 部分遮挡 \\
mixed & 随机组合两种及以上退化 & 复杂复合退化 \\
\bottomrule
\end{tabular}
\end{table}

退化处理在图像输入 CLIP 编码器之前在线完成，无需为每种退化条件预先保存完整的图像副本，从而显著降低存储开销。我们的预期是：真实 EEG 带来的性能增益在退化严重时更为突出，因为此时视觉编码器输出的可靠性降低，模型对 EEG 辅助信息的依赖自然增大。

\subsection{多模态融合机制}

本设计涵盖了从粗粒度到细粒度的多层融合策略。

嵌入（Embedding）级融合将不同模态的向量表示在特征空间中加以组合，常见形式包括直接拼接、加权求和与门控残差融合。本设计采用的门控残差融合公式为 $f_{\mathrm{fused}} = f_{\mathrm{img}} + \sigma(g([f_{\mathrm{img}}; f_{\mathrm{eeg}}])) \odot \Delta(f_{\mathrm{eeg}})$，其中 $g(\cdot)$ 表示门控网络，$\sigma$ 为 sigmoid 激活函数，$\Delta(\cdot)$ 为 EEG 残差变换网络。门控机制使模型能够根据联合输入动态调节 EEG 信号的注入比例：在图像质量较好时弱化 EEG 的影响，而在退化严重时增大 EEG 的修正幅度。

Token 级融合在 CLIP ViT 输出的视觉 patch token 序列中注入 EEG token，通过交叉注意力（cross-attention）实现：视觉 token 充当 Query，EEG token 充当 Key 和 Value，输出经过增强的视觉 token 序列。相比嵌入级融合，token 级融合保留了视觉信息的空间排布结构，与生成式 VLM 的输入格式更为契合。

\begin{table}[h]
\centering
\caption{嵌入级融合与 Token 级融合对比}
\label{tab:fusion-comparison}
\begin{tabular}{llc}
\toprule
特性 & 嵌入级融合 & Token 级融合 \\
\midrule
操作对象 & 全局嵌入向量 & 视觉 patch token 序列 \\
空间信息保留 & 无 & 有 \\
与 VLM 适配性 & 更适合对齐型 & 更适合生成型 \\
\bottomrule
\end{tabular}
\end{table}

\subsection{LoRA 参数高效微调}

大语言模型和视觉语言模型的参数量常达数十亿级别，在小规模训练数据上执行全参数微调极易引发过拟合，并导致模型已有能力的退化。LoRA（Low-Rank Adaptation）通过在注意力层的权重矩阵中插入可训练的低秩分解矩阵来应对这一问题。

给定预训练权重矩阵 $W_0 \in \mathbb{R}^{d\times k}$，LoRA 为其注入低秩增量 $\Delta W = BA$，其中 $B \in \mathbb{R}^{d\times r}$，$A \in \mathbb{R}^{r\times k}$，秩 $r$ 远小于 $d$ 和 $k$。前向计算改写为 $h = W_0 x + \frac{\alpha}{r} \cdot BA x$，$\alpha$ 为缩放因子。训练过程中仅更新 $A$ 与 $B$，原始权重 $W_0$ 保持冻结。

在本设计的生成式描述任务中，LoRA 被应用于 Qwen2-VL 注意力机制的 Query 和 Value 投影矩阵，秩设定为 $r=8$。采用 LoRA 而非全参数微调，是本设计能够在单张 GPU 上完成 Qwen2-VL-2B 模型训练的关键工程前提。
